\begin{textsample}{POS Dim 4 – persona\_gpt – Score 24.00 – t222\_gpt.txt}  \label{ex:f4_pos_029}
When Minister Gwede Mantashe approved Shell ’s seismic \textbf{survey} for oil and gas exploration off South Africa ’s Wild Coast, it came despite strong local opposition and in the immediate aftermath of global climate talks at COP26 that failed to secure the action needed to avert climate breakdown.
Communities, \textbf{fishers}, activists and civil society organisations have been clear : Shell is not welcome on the Wild Coast.
That resistance is rooted not only in concern for the \textbf{ocean} and the climate, but also in Shell ’s track record elsewhere on the African continent.
In Nigeria ’s Niger Delta, Shell ’s \textbf{operations} have been marked by oil \textbf{spills}, severe environmental damage, and decades of protests from communities whose land and \textbf{waters} have been poisoned.
Those protests have often been met with violence, leaving a legacy of trauma and injustice.
For many people along the Wild Coast, this history is a warning of what offshore oil and gas could bring to their own \textbf{shores}.
Seismic blasting is at the heart of the current controversy.
The technique involves firing powerful soundwaves into the \textbf{ocean} to \textbf{map} deposits of oil and gas beneath the \textbf{seabed}.
These blasts threaten marine \textbf{life} in an \textbf{area} known for its rich biodiversity. \textbf{Species} such as southern right \textbf{whales} and ancient coelacanths are part of a fragile \textbf{web} of \textbf{life} that could be disrupted or displaced by intense underwater noise.
The Wild Coast ’s \textbf{ocean} \textbf{ecosystems} also play a vital role in climate \textbf{protection}, and further fossil fuel exploration undermines efforts to keep global warming within safe limits.
For coastal communities, the \textbf{stakes} are immediate and personal.
Many people rely on small-scale \textbf{fishing} and ecotourism for their \textbf{livelihoods}.
Healthy \textbf{fish} populations, intact marine \textbf{ecosystems} and thriving \textbf{wildlife} are the foundation of local economies.
Seismic \textbf{surveys} and the prospect of future \textbf{drilling} raise fears of \textbf{declining} \textbf{fish} \textbf{stocks}, damaged tourism, and the ever-present risk of catastrophic oil \textbf{spills}.
Greenpeace and partner organisations have responded by taking legal action to challenge the approval of Shell ’s \textbf{survey}.
At the same time, protests have spread across the country, uniting a broad coalition of communities, activists and concerned citizens.
On \textbf{beaches}, in towns and cities, people are standing together to defend the Wild Coast and to oppose new fossil fuel projects.
Among the most vocal opponents are communities like Amadiba, who have a long history of resisting destructive projects on their land and along their \textbf{coastline}.
They warn that offshore oil and gas exploration brings the danger of \textbf{spills} and represents yet another form of neo-colonial exploitation, where powerful corporations extract wealth while local people bear the environmental and social costs.
Their call is clear : Shell must be stopped, and decisions about the \textbf{ocean} and the \textbf{coast} must respect the rights, knowledge and needs of the people who live there.
As legal challenges move forward and protests continue, the struggle over Shell ’s plans on the Wild Coast has become a symbol of a wider fight : communities defending their homes, \textbf{oceans} and climate from a fossil fuel industry with a long record of harm.

% matched lemmas: area, beach, coast, coastline, decline, drilling, ecosystem, fish, fisher, fishing, life, livelihood, map, ocean, operation, protection, seabed, shore, specie, spill, stake, stock, survey, water, web, whale, wildlife
\end{textsample}
